\section{One quadratic form for each layer above the Gaussian field}
\label{sec:gaussian}

There is a particularly simple application where a single
unweighted quadratic form tests a whole new layer. Put $K=\mathbb Q(i)$.
For $n\ge1$, let $K_n$ be its ring class field of conductor $2^n$ and set
\[
 Q_n(X,Y)=(2X+Y)^2+2^{2n-2}Y^2\qquad(n\ge2).
\]
For odd rational primes define
$\vartheta_n(x)=\sum_{p\le x,\ p=Q_n(a,b)}\log p$,
counting each represented prime once.

\paragraph{Tower criterion.}
Assume RH for $\zeta_{K_{n-1}}$. Then RH for $\zeta_{K_n}$ is
equivalent to the upper bound
\[
              \vartheta_n(x)\le 2^{-n}x+
                         O_{n,\epsilon}(x^{1/2+\epsilon})
              \quad\text{for every }\epsilon>0 .
\]
In particular these tests for all $n\ge3$ together are equivalent, under
Dirichlet GRH, to RH for the whole tower. Constants may depend on $n$.
The criterion at $n=3$ tests $4X^2+4XY+17Y^2$ with density $1/8$.
Here $K_3=\mathbb Q(i,\sqrt[4]{2})$ (Cox, \S9, ring class field
of $\mathbb Z[\sqrt{-64}]$), and $K_2=\mathbb Q(\zeta_8)$.
Thus the first test counts exactly the quartic nonresidues among
$p\equiv1\pmod8$ considered in the incoming sparse singular-moduli
report, with $2$ as the numerator of the quartic symbol.
The formula $Q_n$ extends this example to every subsequent layer.

\paragraph{Which primes are counted.}
The ring class number formula, including the factor two from the
Gaussian units, gives $[K_n:K]=2^{n-1}$.
The ring class reciprocity law says that an odd prime splits completely
in $K_n$ exactly when $p=a^2+2^{2n}b^2$ for integers $a,b$
(see Cox, \emph{Primes of the form $x^2+ny^2$}, \S\S7,9).
These facts also follow from the exact sequence with quotient
$(\mathbb Z[i]/2^n)^\times/((\mathbb Z/2^n)^\times\mathbb Z[i]^\times)$,
since $\mathbb Z[i]$ has class number one. The fields are nested by
the quotient maps of these groups. For $p\equiv1\pmod 4$ write
$p=a^2+c^2$ with $a$ odd and $c$ even. The pair of absolute values in
this order is unique by unique factorization in $\mathbb Z[i]$.
Thus complete splitting in $K_{n-1}$ and failure of complete splitting
in $K_n$ means $v_2(c)=n-1$. Equivalently $c=2^{n-1}Y$ with $Y$ odd,
and $a=2X+Y$ for some integer $X$. Conversely an odd prime of the displayed
form forces $Y$ odd and has this exact valuation. This proves the promised
description. The quadratic extension $K_2/K$ is Galois over $\mathbb Q$,
and its degree-four Galois group is abelian. Hence Dirichlet GRH supplies
its zeta RH by abelian factorization, starting the induction.

\paragraph{Why an upper bound suffices.}
For any fixed Galois field $M/\mathbb Q$, the series
$P_M(s)=\sum_{p\,\mathrm{split\ completely}}(\log p)p^{-s}$
equals $-[M:\mathbb Q]^{-1}\zeta_M'(s)/\zeta_M(s)$ up to a function
holomorphic on $\Re s>1/2$. Indeed the unramified degree-one prime
ideals occur there in full packets of $[M:\mathbb Q]$; all remaining
unramified terms have exponent at least two. The ramified Euler factors
are harmless in this region. The tested series is $P_{K_{n-1}}-P_{K_n}$
up to any finitely omitted primes. Its pole at 1 has residue
$2^{1-n}-2^{-n}=2^{-n}$. Under RH for the lower field, a zero of
the upper zeta function at $\beta+i\gamma$, $1/2<\beta<1$, contributes
a strictly positive residue equal to its multiplicity divided by $2^n$.
The unsmoothed Landau argument of the preceding sections gives the forward
implication from the bound to zero-freeness. Conversely RH for both
fields and the fixed-field prime ideal explicit formula give the
two-sided $O_n(\sqrt x\log^2(2x))$ error. This proof requires no
holomorphy conjecture for individual Artin factors.

The missing estimate for $Q_3$ illustrates the obstacle to carrying
over the automorphic part of the Lean proof. Detecting $v_2(c)=2$
in a representation $p=a^2+c^2$ uses splitting in a nonabelian
extension of $\mathbb Q$. Indeed the class group at conductor eight
is cyclic of order four: after choosing the odd coordinate as the real
coordinate, the slope $c/a\bmod8$ is an even residue, and doubling slope
2 under Gaussian multiplication gives slope 4 modulo 8. Complex
conjugation inverts this cyclic group. The resulting degree-eight
dihedral group distinguishes this test from restriction to
Dirichlet residue classes. No estimate for this new weight is provided
by the source's theorem about divisibility of its particular theta jets.
